\section{导读：这期访谈为什么重要}

这期访谈把 2026 年初大模型行业的一组关键词放到同一张桌子上：OpenClaw、Agent、后训练、RL scaling、MiMo-V2、多模型编排、1T 基座模型、卡资源分配、组织平权、环境比经验更重要。嘉宾罗福莉曾在阿里达摩院、DeepSeek 工作，目前负责小米大模型团队，访谈中的价值不只在“某个模型表现如何”，而在于它呈现了一个 AI Lab 如何感知范式变化并重组研究方式。

\begin{importantbox}{本期核心命题}
罗福莉的判断是：大模型战争进入第二幕，从 pre-train 主导的 Chat 时代，转向 post-train 主导的 Agent 时代。接下来竞争不只看基座模型，也看 Agent 框架、RL infra、任务环境、成本速度、组织敏捷性和算力分配。
\end{importantbox}

\begin{knowledgebox}{视觉策略说明}
本视频是固定访谈画面，没有教学 slides、白板或产品演示。按本仓库播客标准，正文不重复插入人物帧；封面用于来源识别，正文用概念图和表格承载技术内容。
\end{knowledgebox}

\subsection{本章小结}

本期应当被读作“Agent 时代 AI Lab 运营手册”的口头材料：它讨论的不只是模型，而是模型、框架、后训练、算力和组织如何一起变化。

